\section{Introduction}

\label{sec:intro}

\paragraph{The Fleet Learning Problem.}
Consider an organization deploying $N$ autonomous agents to accomplish similar
tasks---code generation, API integration, document processing, customer
support---across different projects, tenants, or geographic regions. Each
agent independently encounters similar failure modes, discovers similar
solutions, and builds similar intuitions about what strategies work. Without a
sharing mechanism, every agent must rediscover the same lessons from scratch.
Across a fleet of 100 agents, this represents a $100\times$ redundancy in the
experiential learning process.

The natural response is to share what agents learn. But what exactly do agents
learn, and how can it be shared efficiently, safely, and without violating data
sovereignty?

\paragraph{Why Federated Model Training Does Not Apply.}
Classical federated learning~\cite{mcmahan2017communication} addresses
distributed learning by aggregating model gradients or parameters. This
approach is fundamentally mismatched to the agent setting for three reasons.
First, modern agent fleets typically run \emph{frozen} foundation models via
API---there are no local parameters to update. Second, fine-tuning at the
individual-agent level is economically prohibitive (\$10{,}000--\$100{,}000 per
domain) and technically complex. Third, gradient-based federation exposes model
internals and requires synchronized training runs that are incompatible with
continuously operating production agents.

\paragraph{The Key Insight: Share Experiences, Not Weights.}
Active Semantic Optimization (ASO)~\cite{hanzo2026aso} established that
frozen language models can be adapted at decode time via compressed
\emph{experiential priors}---concise, transferable semantic descriptions of
what strategies work for a class of tasks. An experiential prior is not a
gradient; it is structured knowledge extracted from comparing successful and
unsuccessful agent trajectories. Such priors can be shared, aggregated, and
applied without any model parameter access.

This paper extends the single-agent ASO setting to the fleet setting.
We ask: given $N$ agents each running ASO locally, how can they share their
experiential priors efficiently, securely, and with formal convergence
guarantees?

\paragraph{Our Approach.}
We answer this question by composing three existing protocols with new
fleet-level mechanisms:

\begin{enumerate}[leftmargin=1.5em]
    \item \textbf{DSO (ZIP-001/400)~\cite{hanzo2026dso}} provides the
          decentralized protocol for broadcasting and storing experiential
          priors on a peer-to-peer network backed by content-addressed storage.
    \item \textbf{BitDelta (ZIP-007)~\cite{hanzo2026bitdelta}} provides 1-bit
          quantization of experience embedding deltas, achieving ${\sim}10\times$
          compression with bounded reconstruction error, making fleet-scale
          broadcasting practical on commodity networks.
    \item \textbf{DeltaSoup} is our new Byzantine-robust aggregation layer
          that applies geometric median filtering to received compressed
          experiences before merging them into the local library, tolerating
          up to $f < N/3$ adversarial contributors.
\end{enumerate}

\paragraph{Contributions.}
This paper makes the following contributions:
\begin{enumerate}[leftmargin=1.5em]
    \item \textbf{Agent-level DSO integration}: A complete protocol for
          extracting, compressing, and broadcasting agent experiences via the
          DSO network, including the \texttt{CompressedBatch} wire format.
    \item \textbf{BitDelta experience compression}: Application of 1-bit delta
          quantization to experience embeddings, with analysis showing
          $29.5\times$ theoretical and $10\times$ practical compression ratios
          and cosine similarity preservation bounds.
    \item \textbf{DeltaSoup aggregation}: A Byzantine-robust aggregation
          protocol for compressed experience batches using coordinate-wise
          weighted median, with formal $(f,\varepsilon)$-robustness guarantees.
    \item \textbf{Fleet convergence analysis}: Proof that the FAI protocol
          converges monotonically under honest majority with rate depending on
          network connectivity and experience diversity.
    \item \textbf{Empirical evaluation}: Simulation results with $N=10, 50, 100$
          agents showing fleet improvement over isolated agents and resilience to
          Byzantine adversaries.
\end{enumerate}

\paragraph{Paper Outline.}
Section~\ref{sec:background} reviews background on federated learning, DSO, BitDelta,
and DeltaSoup. Section~\ref{sec:architecture} describes the system architecture.
Section~\ref{sec:bitdelta} formalizes BitDelta experience compression.
Section~\ref{sec:deltasoup} develops DeltaSoup aggregation with formal guarantees.
Section~\ref{sec:convergence} analyzes fleet learning dynamics and convergence.
Section~\ref{sec:harness} describes integration with the Hanzo agent harness.
Section~\ref{sec:privacy} addresses privacy and security.
Section~\ref{sec:evaluation} presents experimental evaluation.
Section~\ref{sec:discussion} discusses future directions.
Section~\ref{sec:related} reviews related work.
Section~\ref{sec:conclusion} concludes.

% ============================================================================
% 2. BACKGROUND
% ============================================================================
